\documentclass{ximera}

\input{../../../preamble.tex}


\definecolor{fvdnavy}{HTML}{071D43}
\definecolor{fvdcyan}{HTML}{20BFC6}
\definecolor{fvdcyanlight}{HTML}{EFFCFB}
\definecolor{fvdmagenta}{HTML}{F10D69}
\definecolor{fvdmagentalight}{HTML}{FFF1F7}
\definecolor{fvdtext}{HTML}{163044}





%=================================================
% Pedagogische omgevingen
% PDF: tcolorbox
% HTML: learning-box
%=================================================

\newenvironment{voorkennis}
{%
    \ifdefined\HCode
        \HCode{%
<section class="learning-box learning-box--prior">
<div class="learning-box__header">
<span class="learning-box__icon" aria-hidden="true"></span>
<span class="learning-box__title">Voorkennis</span>
</div>
<div class="learning-box__content">
        }%
        \begin{itemize}
    \else
       \begin{tcolorbox}[
    enhanced,
    breakable,
    colback=fvdcyanlight,
    colframe=fvdcyan,
    coltext=fvdtext,
    boxrule=0.5pt,
    leftrule=4pt,
    arc=4mm,
    outer arc=4mm,
    left=7mm,
    right=6mm,
    top=5mm,
    bottom=5mm,
    before skip=6mm,
    after skip=6mm,
    title=Voorkennis,
    coltitle=fvdcyan,
    fonttitle=\bfseries\Large,
    attach title to upper={\par\smallskip},
    boxed title style={
        empty,
        left=0mm,
        right=0mm,
        top=0mm,
        bottom=0mm
    },
    overlay={
        \draw[
            fvdcyan,
            line width=0.5pt
        ]
        ([xshift=42mm,yshift=-13mm]frame.north west)
        --
        ([xshift=-7mm,yshift=-13mm]frame.north east);
    }
]
        \begin{itemize}
    \fi
}
{%
    \end{itemize}
    \ifdefined\HCode
        \HCode{%
</div>
</section>
        }%
    \else
        \end{tcolorbox}
    \fi
}


\newenvironment{leerdoelen}
{%
    \ifdefined\HCode
        \HCode{%
<section class="learning-box learning-box--goals">
<div class="learning-box__header">
<span class="learning-box__icon" aria-hidden="true"></span>
<span class="learning-box__title">Leerdoelen</span>
</div>
<div class="learning-box__content">
        }%
        \begin{itemize}
    \else
\begin{tcolorbox}[
    enhanced,
    breakable,
    colback=fvdmagentalight,
    colframe=fvdmagenta,
    coltext=fvdtext,
    boxrule=0.5pt,
    leftrule=4pt,
    arc=4mm,
    outer arc=4mm,
    left=7mm,
    right=6mm,
    top=5mm,
    bottom=5mm,
    before skip=6mm,
    after skip=6mm,
    title=Leerdoelen,
    coltitle=fvdmagenta,
    fonttitle=\bfseries\Large,
    attach title to upper={\par\smallskip},
    boxed title style={
        empty,
        left=0mm,
        right=0mm,
        top=0mm,
        bottom=0mm
    },
    overlay={
        \draw[
            fvdmagenta,
            line width=0.5pt
        ]
        ([xshift=42mm,yshift=-13mm]frame.north west)
        --
        ([xshift=-7mm,yshift=-13mm]frame.north east);
    }
]
        \begin{itemize}
    \fi
}
{%
    \end{itemize}
    \ifdefined\HCode
        \HCode{%
</div>
</section>
        }%
    \else
        \end{tcolorbox}
    \fi
}


\addPrintStyle{../../..}

\title{Van werkelijkheid naar wiskunde --- en terug}
\author{Ingmar Herreman}



\begin{document}








\maketitle


% ====================================================================
% 1. GROOTHEDEN EN VARIABELEN
% ====================================================================

\begin{oefening}[Van experiment naar variabelen]{\oefB \ESS}

Bij een experiment wordt onderzocht hoe de uitrekking van een veer
afhangt van de kracht die erop wordt uitgeoefend.

Voor verschillende krachten wordt de uitrekking gemeten.

\begin{question}
Welke grootheid is de onafhankelijke variabele?

\begin{multipleChoice}
  \choice[correct]{de uitgeoefende kracht}
  \choice{de uitrekking van de veer}
\end{multipleChoice}


\begin{oplossing}
De onafhankelijke variabele is de \textbf{uitgeoefende kracht}.
\end{oplossing}

\end{question}

\begin{question}
Welke grootheid is de afhankelijke variabele?

\begin{multipleChoice}
  \choice{de uitgeoefende kracht}
  \choice[correct]{de uitrekking van de veer}
\end{multipleChoice}


\begin{oplossing}
De afhankelijke variabele is de \textbf{uitrekking van de veer}.
\end{oplossing}

\end{question}

\begin{question}
Welke keuze van variabelen is geschikt?

\begin{multipleChoice}
  \choice[correct]{$F=$ kracht en $u=$ uitrekking}
  \choice{$F=$ kleur van de veer en $u=$ temperatuur van het lokaal}
  \choice{$F=$ uitrekking en $u=$ merk van de veer}
\end{multipleChoice}


\begin{oplossing}
Een geschikte keuze is
\[
F=\text{kracht},
\qquad
u=\text{uitrekking}.
\]
\end{oplossing}

\end{question}

\end{oefening}


% ====================================================================
% 2. TABEL EN GRAFIEK INTERPRETEREN
% ====================================================================

\begin{oefening}[Een afkoelend voorwerp]{\oefB \ESS}

Van een warm voorwerp wordt om de vijf minuten de temperatuur gemeten.

\[
\begin{array}{c|cccccc}
t\;(\mathrm{min})
&0&5&10&15&20&25\\ \hline
T\;({}^{\circ}\mathrm C)
&82&67&56&48&42&38
\end{array}
\]

\begin{question}
Welke uitspraak beschrijft de gegevens het best?

\begin{multipleChoice}
  \choice{De temperatuur stijgt ongeveer gelijkmatig.}
  \choice{De temperatuur daalt elke vijf minuten met exact hetzelfde bedrag.}
  \choice[correct]{De temperatuur daalt, maar de daling wordt geleidelijk kleiner.}
  \choice{Er is geen herkenbaar patroon.}
\end{multipleChoice}


\begin{oplossing}
De temperatuur daalt, maar de daling wordt geleidelijk kleiner.
\end{oplossing}

\end{question}

\begin{question}
Je wilt de gemeten temperatuur op $t=15\,\mathrm{min}$ kennen.
Welke voorstelling geeft die exacte meetwaarde het gemakkelijkst?

\begin{multipleChoice}
  \choice[correct]{de tabel}
  \choice{een schets van de grafiek}
\end{multipleChoice}


\begin{oplossing}
De tabel geeft de exacte gemeten waarde het gemakkelijkst:
\[
T(15)=48\,{}^\circ\mathrm C.
\]
\end{oplossing}

\end{question}

\begin{question}
Je wilt vooral het globale verloop van de temperatuur onderzoeken.
Welke voorstelling is dan meestal het handigst?

\begin{multipleChoice}
  \choice{de tabel}
  \choice[correct]{de grafiek}
\end{multipleChoice}


\begin{oplossing}
Voor het globale verloop is een \textbf{grafiek} doorgaans het handigst.
\end{oplossing}

\end{question}

\end{oefening}


% ====================================================================
% 3. MATHEMATISEREN EN DEMATHEMATISEREN
% ====================================================================

\begin{oefening}[Een batterij ontlaadt]{\oefB \ESS}

Tijdens een experiment wordt de spanning van een batterij gedurende
enkele uren gemeten.

Een leerling kiest

\[
t=\text{tijd in uren},
\qquad
U=\text{spanning in volt}.
\]

Een model geeft

\[
U(4)=10{,}8.
\]

\begin{question}
Welke stap is een voorbeeld van mathematiseren?

\begin{multipleChoice}
  \choice[correct]{De tijd en spanning kiezen als relevante grootheden en ze voorstellen door $t$ en $U$.}
  \choice{Besluiten dat de spanning na vier uur volgens het model $10{,}8$ volt bedraagt.}
\end{multipleChoice}


\begin{oplossing}
Het kiezen van de relevante grootheden en ze voorstellen door $t$ en $U$
is een stap van \textbf{mathematiseren}.
\end{oplossing}

\end{question}

\begin{question}
Welke interpretatie van $U(4)=10{,}8$ is correct?

\begin{multipleChoice}
  \choice{Na $10{,}8$ uur bedraagt de spanning $4$ volt.}
  \choice[correct]{Na $4$ uur bedraagt de spanning volgens het model $10{,}8$ volt.}
  \choice{De spanning daalt elk uur met $10{,}8$ volt.}
\end{multipleChoice}


\begin{oplossing}
\[
U(4)=10{,}8
\]
betekent: na $4$ uur bedraagt de spanning volgens het model
$10{,}8\,\mathrm V$.
\end{oplossing}

\end{question}

\end{oefening}


% ====================================================================
% 4. INTERPOLATIE EN EXTRAPOLATIE
% ====================================================================

\begin{oefening}[Een groeiende bacteriecultuur]{\oefB \ESS}

Een bacteriecultuur wordt gedurende $12$ uur gevolgd.
Er zijn metingen beschikbaar voor

\[
0\leq t\leq12.
\]

Een model wordt opgesteld op basis van deze gegevens.

\begin{question}
Een voorspelling voor $t=7{,}5$ uur is

\begin{multipleChoice}
  \choice[correct]{interpolatie.}
  \choice{extrapolatie.}
\end{multipleChoice}


\begin{oplossing}
Omdat $7{,}5$ binnen het meetinterval $[0,12]$ ligt, is dit
\textbf{interpolatie}.
\end{oplossing}

\end{question}

\begin{question}
Een voorspelling voor $t=15$ uur is

\begin{multipleChoice}
  \choice{interpolatie.}
  \choice[correct]{extrapolatie.}
\end{multipleChoice}


\begin{oplossing}
Omdat $15$ buiten het meetinterval $[0,12]$ ligt, is dit
\textbf{extrapolatie}.
\end{oplossing}

\end{question}

\begin{question}
Een leerling gebruikt het model om de populatie na $200$ uur te voorspellen.
Welke beoordeling is het best?

\begin{multipleChoice}
  \choice{De voorspelling is betrouwbaar zolang de berekening correct is.}
  \choice[correct]{We extrapoleren zeer ver en moeten nagaan of de aannames en het patroon dan nog geldig zijn.}
  \choice{Een wiskundig model mag nooit buiten de meetpunten gebruikt worden.}
\end{multipleChoice}


\begin{oplossing}
Dit is een zeer verre extrapolatie. We moeten nagaan of de aannames
van het model en het waargenomen patroon na zoveel tijd nog geldig zijn.
Een correcte berekening alleen maakt de voorspelling niet betrouwbaar.
\end{oplossing}

\end{question}

\end{oefening}


% ====================================================================
% 5. EEN MODEL KRITISCH BEOORDELEN
% ====================================================================

\begin{oefening}[Vrije val: wat zit niet in het model?]{\oefU \ESS}

Voor een vallend voorwerp wordt het model

\[
h(t)=100-4{,}9t^2
\]

gebruikt, met $h$ in meter en $t$ in seconden.

\begin{question}
Welke informatie is in dit model rechtstreeks opgenomen?

\begin{multipleChoice}
  \choice[correct]{de hoogte als functie van de tijd}
  \choice{de vorm en kleur van het voorwerp}
  \choice{de windsnelheid op elke hoogte}
  \choice{alle eigenschappen van de werkelijke valbeweging}
\end{multipleChoice}


\begin{oplossing}
Het model beschrijft rechtstreeks de \textbf{hoogte als functie van de tijd}.
\end{oplossing}

\end{question}

\begin{question}
Het model voorspelt voor een bepaalde tijd een negatieve hoogte.
Wat is de beste reactie?

\begin{multipleChoice}
  \choice{Een negatieve uitkomst bewijst dat de algebra fout is.}
  \choice[correct]{We moeten nagaan of het model nog gebruikt wordt binnen de fysische situatie waarvoor het bedoeld was.}
  \choice{Negatieve getallen mogen niet in modellen voorkomen.}
\end{multipleChoice}


\begin{oplossing}
Een negatieve hoogte betekent niet automatisch dat de algebra fout is.
We moeten nagaan of het model nog gebruikt wordt binnen de fysische
situatie waarvoor het bedoeld was. Na het bereiken van de grond
beschrijft dit model de werkelijke valbeweging bijvoorbeeld niet meer.
\end{oplossing}

\end{question}

\end{oefening}


% ====================================================================
% 6. MISLEIDENDE GRAFIEKEN
% ====================================================================

\begin{oefening}[Dezelfde gegevens, een andere indruk]{\oefU \ESS}

De gemiddelde meetwaarde van een grootheid stijgt van $198$ naar $202$.

Een eerste grafiek gebruikt een verticale as van $0$ tot $220$.
Een tweede grafiek gebruikt een verticale as van $197$ tot $203$.

\begin{question}
In welke grafiek zal de stijging visueel het grootst lijken?

\begin{multipleChoice}
  \choice{in de eerste grafiek}
  \choice[correct]{in de tweede grafiek}
  \choice{in beide noodzakelijk even groot}
\end{multipleChoice}


\begin{oplossing}
De stijging zal in de tweede grafiek visueel het grootst lijken,
omdat de verticale as daar slechts van $197$ tot $203$ loopt.
\end{oplossing}

\end{question}

\begin{question}
Is de tweede grafiek daarom automatisch fout?

\begin{multipleChoice}
  \choice{Ja, iedere verticale as moet bij nul beginnen.}
  \choice[correct]{Nee, maar de gekozen schaal moet mee in rekening worden gebracht bij de interpretatie.}
\end{multipleChoice}


\begin{oplossing}
Nee. De grafiek is niet automatisch fout, maar de gekozen schaal
vergroot de visuele indruk van het verschil. De schaal moet dus
mee geïnterpreteerd worden.
\end{oplossing}

\end{question}

\begin{question}
Welke informatie controleer je minstens voordat je een conclusie trekt?

\begin{selectAll}
  \choice[correct]{de grootheden op de assen}
  \choice[correct]{de gebruikte eenheden}
  \choice[correct]{de schaal van de assen}
  \choice[correct]{het bereik dat wordt weergegeven}
  \choice{of de grafiek er visueel aantrekkelijk uitziet}
\end{selectAll}


\begin{oplossing}
Controleer minstens:
\begin{itemize}
  \item de grootheden op de assen;
  \item de gebruikte eenheden;
  \item de schaal van de assen;
  \item het weergegeven bereik.
\end{itemize}
\end{oplossing}

\end{question}

\end{oefening}


% ====================================================================
% 7. SAMENHANG EN CAUSALITEIT
% ====================================================================

\begin{oefening}[Zonnebrillen en ijsjes]{\oefU \ESS}

In een onderzoek over verschillende dagen blijkt dat op dagen waarop
meer zonnebrillen worden verkocht, ook meer ijsjes worden verkocht.

Een leerling concludeert:

\begin{center}
``De verkoop van zonnebrillen veroorzaakt een hogere ijsverkoop.''
\end{center}

\begin{question}
Is die conclusie verantwoord?

\begin{multipleChoice}
  \choice{Ja, want beide variabelen veranderen samen.}
  \choice[correct]{Nee, samenhang alleen bewijst geen causaliteit.}
\end{multipleChoice}


\begin{oplossing}
Nee. Dat beide variabelen samen veranderen toont \textbf{samenhang},
maar bewijst op zichzelf geen causaliteit.
\end{oplossing}

\end{question}

\begin{question}
Welke derde variabele zou de gevonden samenhang kunnen verklaren?

\begin{multipleChoice}
  \choice[correct]{de buitentemperatuur}
  \choice{het serienummer van de zonnebrillen}
  \choice{de volgorde waarin de gegevens in de tabel staan}
\end{multipleChoice}


\begin{oplossing}
Een mogelijke derde variabele is de \textbf{buitentemperatuur}.
Op warme dagen kunnen zowel meer zonnebrillen als meer ijsjes worden verkocht.
\end{oplossing}

\end{question}

\end{oefening}


% ====================================================================
% 8. INTEGRERENDE MODELOPGAVE
% ====================================================================

\begin{oefening}[Een zonnepaneel onderzoeken]{\oefU \ESS}

Tijdens één dag wordt op verschillende tijdstippen het elektrisch
vermogen van een zonnepaneel gemeten.

\[
\begin{array}{c|ccccccc}
t\;(\mathrm h)
&8&10&12&14&16&18&20\\ \hline
P\;(\mathrm W)
&85&210&345&390&315&155&25
\end{array}
\]

\begin{question}
Welke twee grootheden worden in dit experiment onderzocht?

\begin{freeResponse}
\end{freeResponse}


\begin{oplossing}
De onderzochte grootheden zijn:
\begin{itemize}
  \item de tijd $t$, uitgedrukt in uur;
  \item het elektrisch vermogen $P$, uitgedrukt in watt.
\end{itemize}
\end{oplossing}

\end{question}

\begin{question}
Welke variabele zou je als onafhankelijke variabele kiezen? Verklaar.

\begin{freeResponse}
\end{freeResponse}


\begin{oplossing}
Een natuurlijke keuze is de tijd $t$ als onafhankelijke variabele,
omdat onderzocht wordt hoe het vermogen gedurende de dag verandert.
Het vermogen $P$ is dan de afhankelijke variabele.
\end{oplossing}

\end{question}

\begin{question}
Beschrijf het globale patroon in de gegevens.

\begin{freeResponse}
\end{freeResponse}


\begin{oplossing}
Het vermogen stijgt van $8$ uur tot ongeveer $14$ uur, bereikt daar
een maximum en daalt daarna opnieuw. Het globale patroon is dus
duidelijk \textbf{niet lineair}.
\end{oplossing}

\end{question}

\begin{question}
Is een schatting van $P$ om $13$ uur interpolatie of extrapolatie?
Verklaar.

\begin{freeResponse}
\end{freeResponse}


\begin{oplossing}
Dit is \textbf{interpolatie}, want $13$ uur ligt binnen het
meetinterval van $8$ uur tot $20$ uur.
\end{oplossing}

\end{question}

\begin{question}
Een leerling gebruikt deze gegevens om het vermogen om $2$ uur 's nachts
te voorspellen door het patroon eenvoudig verder te trekken.

Waarom moet je die voorspelling kritisch beoordelen?

\begin{freeResponse}
\end{freeResponse}


\begin{oplossing}
$2$ uur 's nachts ligt ver buiten het gemeten tijdsinterval.
Bovendien verandert de fysische situatie: 's nachts is er geen
rechtstreekse zoninstraling. Het patroon van overdag mag dus niet
zomaar verder worden doorgetrokken.
\end{oplossing}

\end{question}

\begin{question}
Noem minstens één factor uit de werkelijkheid die het vermogen beïnvloedt
maar niet rechtstreeks in een model $P(t)$ voorkomt.

\begin{freeResponse}
\end{freeResponse}


\begin{oplossing}
Mogelijke factoren zijn bijvoorbeeld:
\begin{itemize}
  \item bewolking;
  \item zoninstraling;
  \item schaduw;
  \item de oriëntatie en hellingshoek van het paneel;
  \item de temperatuur van het paneel.
\end{itemize}
Andere wetenschappelijk verantwoorde antwoorden zijn mogelijk.
\end{oplossing}

\end{question}

\end{oefening}


% ====================================================================
% 9. OPEN MODELOPGAVE
% ====================================================================

\begin{oefening}[Van werkelijkheid naar model]{\oefG}

Je wilt onderzoeken hoe snel een kop warme drank afkoelt in een lokaal.

Ontwerp een eenvoudige aanpak voor dit onderzoek.

Beschrijf minstens:

\begin{enumerate}
  \item welke grootheden je meet;
  \item welke variabelen je invoert en welke eenheden je gebruikt;
  \item welke grootheid je als onafhankelijke variabele kiest;
  \item hoe je de gegevens zou voorstellen;
  \item welke patronen je in de gegevens zou onderzoeken;
  \item welke aannames of omstandigheden belangrijk kunnen zijn;
  \item wanneer een voorspelling met je model volgens jou onbetrouwbaar
        zou kunnen worden.
\end{enumerate}



\begin{oplossing}
Een mogelijke onderzoeksopzet is:

\begin{enumerate}
  \item Meet de \textbf{tijd} en de \textbf{temperatuur van de drank}
        op regelmatige tijdstippen. Het is ook zinvol de
        \textbf{omgevingstemperatuur} te registreren.

  \item Kies bijvoorbeeld
        \[
        t=\text{tijd in minuten},
        \qquad
        T=\text{temperatuur in }{}^\circ\mathrm C.
        \]

  \item De tijd $t$ is de onafhankelijke variabele en de temperatuur
        $T$ de afhankelijke variabele.

  \item Noteer de meetwaarden in een tabel en stel ze vervolgens
        grafisch voor.

  \item Onderzoek hoe de temperatuur afneemt en of de afkoeling
        geleidelijk trager wordt naarmate de temperatuur van de drank
        dichter bij de omgevingstemperatuur komt.

  \item Houd relevante omstandigheden zoveel mogelijk constant,
        bijvoorbeeld de hoeveelheid drank, het soort beker, een
        eventueel deksel, luchtstroming en de omgevingstemperatuur.

  \item Voorspellingen worden minder betrouwbaar wanneer we ver buiten
        het gemeten tijdsinterval extrapoleren of wanneer de
        omstandigheden veranderen.
\end{enumerate}

Dit is één mogelijke aanpak. Andere goed gemotiveerde en
wetenschappelijk verantwoorde onderzoeksopzetten zijn eveneens correct.
\end{oplossing}

\end{oefening}

\end{document}